\documentclass[10pt,a4paper]{amsart}
\usepackage[margin=19mm]{geometry}
\usepackage{amsmath,amssymb,mathtools}
\usepackage[scr=rsfs,cal=euler]{mathalfa}
\usepackage{xfrac}
\usepackage[unicode,hidelinks,bookmarks=false]{hyperref}
\newcommand{\ie}{i.e.\ }
\newcommand{\cf}{cf.\ }
\newcommand{\resp}{respectively\ }
\newcommand{\Subsection}{\subsection}
\providecommand{\nobreakdash}{\nobreak\hbox{-}}
\setlength{\emergencystretch}{2em}
\multlinegap0pt
\usepackage{lmodern}
\numberwithin{equation}{section}
\newtheorem{proposition}{Proposition}
\newtheorem{corollary}{Corollary}
\newtheorem{definition}{Definition}
\newtheorem{example}{Example}
\newtheorem{remark}{Remark}
\newcommand{\id}{\operatorname{id}}
\newcommand{\tr}{\operatorname{tr}}
\begin{document}
\title[Contact canonoid maps and their conserved and dissipated qua]{Contact canonoid maps and their conserved and dissipated quantities}
\author{R. Azuaje}
\date{}
\maketitle
\noindent\emph{Source and licence.} Contact canonoid maps and their conserved and dissipated quantities. Version arXiv:2609.18661v1 (CC BY 4.0). This material is distributed under the Creative Commons Attribution 4.0 International licence, http://creativecommons.org/licenses/by/4.0/, as recorded in the arXiv OAI licence metadata for this exact version and shown on the abstract page https://arxiv.org/abs/2609.18661v1. Full rights notice: licenses/arxiv-2609.18661-CC-BY-4.0.txt.\par\smallskip
\noindent\textit{Editorial scope.} Counted unit: the original \texttt{proposition} environment \texttt{prop:tensor} at source lines 437--490 together with its complete proof at lines 492--596; the delivered document keeps source lines 414--597 so that every referenced label and every citation resolves inside it. Native editable source is the paper's own concatenated TeX; the AMS wrapper is editorial formatting only. No publisher class, upstream script or unrelated artwork is redistributed.
\section{An invariant tensor and its polynomial invariants}
\label{sec:tensor}

We now develop a tensorial counterpart of the
Hamiltonian-normalized polynomial construction.
For a contact canonoid map, the pulled-back form
$\alpha=\Phi^*\eta$ and the associated function
$K=-\alpha(X_H)$ determine a smooth invariant endomorphism
on $\{K\neq0\}$. Its traces are conserved quantities,
and their products with $H$ are dissipated quantities.
The construction accommodates singular maps and remains
smooth at zeros of $H$, where the endomorphism vanishes.
We establish the precise relation between its traces and
the coefficients of the Hamiltonian-normalized polynomial,
showing that these constructions encode the same spectral
information and yield at most $n$ functionally independent
trace invariants in dimension $2n+1$. We conclude with an
explicit formula in dimension three.



\subsection{The intrinsic tensor}


\begin{proposition}
\label{prop:tensor}
Let $\Phi$ be a contact canonoid map for $(M,\eta,H)$, and set
\[
X=X_H,\qquad a=R(H),\qquad
\alpha=\Phi^*\eta,\qquad K=-\alpha(X).
\]
On the open set
\[
V=\{x\in M:K(x)\neq0\},
\]
define
\[
\sigma=-\frac{\alpha}{K},
\qquad
\mathcal D=\ker(\eta|_V).
\]
There exists a unique smooth endomorphism $N:TV\to TV$ such that
\[
\eta(NY)=0,
\qquad
d\eta(NY,Z)=-H\,d\sigma(Y,Z)
\]
for every vector field $Y$ on $V$ and every section $Z$ of
$\mathcal D$. Moreover,
\[
NX=0,\qquad \mathcal L_XN=0,
\]
and $N_x=0$ at every point $x\in V$ for which $H(x)=0$.

On the open subset
\[
U=\{x\in M:H(x)K(x)\neq0\},
\]
put $\lambda=-\eta/H$. Then $N$ is equivalently characterized by
\[
\lambda(NY)=0,
\qquad
d\lambda(NY,Z)=d\sigma(Y,Z)
\]
for all vector fields $Y,Z$ on $U$.

For every integer $k\geq1$, the functions
\[
I_k=\operatorname{tr}(N^k),
\qquad
D_k=H\operatorname{tr}(N^k)
\]
are smooth on $V$ and satisfy
\[
X(I_k)=0,\qquad X(D_k)=-aD_k.
\]
No invertibility of $\Phi$ is assumed.
\end{proposition}

\begin{proof}
The canonoid condition and the balance law for $K$ give
\[
\mathcal L_X\alpha=-b\alpha,
\qquad
X(K)=-bK.
\]
Consequently,
\[
\mathcal L_X\sigma=0,
\qquad
\sigma(X)=1.
\]
Cartan's formula therefore yields
\[
\iota_Xd\sigma=0.
\]

Since $\eta$ is contact, the restriction of $d\eta$ to
$\mathcal D$ is nondegenerate. Thus, for each vector field $Y$,
there is a unique section $NY$ of $\mathcal D$ satisfying
\[
d\eta(NY,Z)=-H\,d\sigma(Y,Z),
\qquad Z\in\Gamma(\mathcal D).
\]
The inverse of the corresponding bundle isomorphism
$\mathcal D\to\mathcal D^*$ is smooth, so this construction
defines a smooth endomorphism $N$. Taking $Y=X$ gives $NX=0$.
At a point where $H=0$, nondegeneracy gives $N=0$.

We next prove invariance. Since
\[
\mathcal L_X\eta=-a\eta,
\]
the distribution $\mathcal D$ is invariant under $X$.
Lie differentiation of $\eta(NY)=0$ gives
\[
\eta\bigl((\mathcal L_XN)Y\bigr)=0.
\]
Furthermore,
\[
\mathcal L_Xd\eta=-da\wedge\eta-a\,d\eta,
\qquad
\mathcal L_Xd\sigma=0,
\qquad
X(H)=-aH.
\]
We may therefore Lie differentiate the defining identity for
$N$, with its second argument restricted to $\mathcal D$.
For $Z\in\Gamma(\mathcal D)$ this gives
\[
(\mathcal L_Xd\eta)(NY,Z)
+d\eta\bigl((\mathcal L_XN)Y,Z\bigr)
=aH\,d\sigma(Y,Z).
\]
Both $NY$ and $Z$ lie in $\mathcal D$, so
\[
(\mathcal L_Xd\eta)(NY,Z)
=-a\,d\eta(NY,Z)
=aH\,d\sigma(Y,Z).
\]
Hence
\[
d\eta\bigl((\mathcal L_XN)Y,Z\bigr)=0
\qquad\text{for every }Z\in\Gamma(\mathcal D).
\]
Together with
$\eta((\mathcal L_XN)Y)=0$, nondegeneracy on $\mathcal D$
implies $\mathcal L_XN=0$.

On $U$, the form $\lambda=-\eta/H$ satisfies
\[
\lambda(X)=1,\qquad
\mathcal L_X\lambda=0,\qquad
\iota_Xd\lambda=0.
\]
Since $\ker\lambda=\mathcal D|_U$, we have
\[
TU=\mathcal D|_U\oplus\operatorname{span}\{X\}.
\]
For $Z\in\Gamma(\mathcal D|_U)$, the identity
\[
d\lambda=-\frac{1}{H}d\eta
+\frac{1}{H^2}dH\wedge\eta
\]
and the condition $\eta(NY)=0$ give
\[
d\lambda(NY,Z)
=-\frac{1}{H}d\eta(NY,Z)
=d\sigma(Y,Z).
\]
Both sides vanish when $Z=X$, so the same identity holds
for arbitrary $Z$. Conversely, restricting this identity
to $Z\in\Gamma(\mathcal D|_U)$ recovers the defining
equation for $N$.

Finally, Lie differentiation is a derivation for composition
and commutes with contraction. Therefore,
\[
X\bigl(\operatorname{tr}(N^k)\bigr)
=\operatorname{tr}\bigl(\mathcal L_X(N^k)\bigr)=0.
\]
Multiplication by $H$, together with $X(H)=-aH$, proves
the balance law for $D_k$.
\end{proof}

\end{document}
